\section{Introduction}

Intelligent systems have been developed for years to model, predict, and
simulate complex processes. Neuroscience is one of the most demanding domains
for these systems because the information is heterogeneous, dynamic, and
distributed across several biological scales. In the particular case of the
mouse brain, researchers can currently access public resources that include
large-scale electrophysiological recordings, anatomical atlases, visual
prediction benchmarks, and local connectomic reconstructions. This fact opens
the possibility of building digital models that are more realistic than previous
simplified simulations.

Nevertheless, this new scenario introduces a relevant problem. The availability
of more data does not automatically imply that a model can be interpreted as a
mechanistic explanation of the brain. A model may predict neural responses with
good performance and, at the same time, fail to identify the anatomical
structure producing those responses. Similarly, a reproducible experimental
signal can be scientifically useful without validating a specific connectivity
mechanism. This issue becomes more critical when the expression ``digital twin''
is used, since it can suggest a complete and causally faithful replica when the
available data usually cover only partial aspects of the biological system.

In this paper, the MouseBrainBench framework is proposed. It is a reproducible
software and experimental architecture designed to audit partial digital models
of the mouse brain. The aim is not to build a complete digital mouse brain, nor
to replace consolidated simulators and public resources. The objective is to
provide a critical layer able to integrate these resources, evaluate the
evidence produced by models, and determine which scientific claims are actually
supported.

The proposal is based on the separation of evidence. Predictive accuracy,
repeat reliability, anatomical specificity, directed identifiability,
structure-function association, perturbation behaviour, and computational cost
are represented as different blocks. Thus, the framework avoids promoting a
single positive metric to a broad mechanistic statement. This conservative
design is especially important in the context of high-impact Artificial
Intelligence and neuroscience research, where methodological transparency and
claim control are essential.

Several experiments have been designed to validate the proposal. Synthetic
known-truth cases are used to test the logic of the Mechanistic Identifiability
Score. Allen Visual Behavior Neuropixels is used as a negative mechanistic
control, since reproducible responses are detected but anatomical and directed
identifiability gates do not pass. Sensorium static and Dynamic Sensorium are
used as predictive benchmarks, showing the difference between prediction and
mechanistic interpretation. Finally, MICRONS CAVE data are used as the main
empirical positive case, where local structure-function associations can be
evaluated under anatomical controls.

The rest of the article is structured as follows. Section~2 details the
contributions of the proposal. Section~3 introduces the foundations and related
work. Section~4 presents the MouseBrainBench framework and its components.
Section~5 describes the datasets and their specific role in the proposal.
Section~6 focuses on the experiments carried out to validate the system.
Finally, Sections~7--10 discuss the results, limitations, methodological
details, and reproducibility information.
